\chapter{Complexity — Engine}
%%% generated by the blueprint pipeline from TCSlib.Complexity.SpaceComplexity.Codec.Engine
\section{Overview}

Declarations of the Lean module \texttt{TCSlib.Complexity.SpaceComplexity.Codec.Engine}.

\section{Declarations}

\begin{theorem}[A Chain of Register Increments]
\lean{Complexity.Codec.incChain}
\label{Complexity.Codec.incChain}
\leanok
\uses{Complexity.LogProg.ARM, Complexity.LogProg.AReach, Complexity.LogProg.GuardLe, Complexity.LogProg.arm_chain}
\difficulty{1}
Let $A$ be a register program with $m$ registers and $d$ Boolean-oracle queries, let $y$
be its input, let $r$ be a register, and let $(\ell_i)_{i\in\mathbb{N}}$ be a sequence of
labels. Suppose that for every $i<k$, the instruction at label $\ell_i$ increments
register $r$ and proceeds to $\ell_{i+1}$. If the initial register assignment $v$ has
all values at most $B$ and $v(r)+k\le B$, then the program reaches $\ell_k$ from
$\ell_0$, with register $r$ increased by $k$ and every other register unchanged.
Throughout the run the prescribed safety condition holds and every register remains
at most $B$.
\end{theorem}

\begin{theorem}[Repeated decrement execution]
\lean{Complexity.Codec.decChain}
\label{Complexity.Codec.decChain}
\leanok
\uses{Complexity.LogProg.ARM, Complexity.LogProg.AReach, Complexity.LogProg.GuardLe, Complexity.LogProg.arm_chain}
\difficulty{1}
Let $A$ be a register program with $m$ registers and $d$ oracle queries, let $\Lambda$
be its set of labels, and fix an input $y$, a register $r$, labels $\ell_i$ for
$i\in\mathbb N$, and natural numbers $k,B$. Suppose that for every $i<k$, label $\ell_i$
instructs the program to decrement register $r$ and continue at $\ell_{i+1}$; initially
every register is at most $B$ and the value in register $r$ is at least $k$. Then the
program has a finite execution from label $\ell_0$ to label $\ell_k$, with register $r$
reduced by $k$ and all other register values unchanged. Every configuration before the
endpoint satisfies the program's guard: its instruction is safe to execute and every
register remains at most $B$.
\end{theorem}

\begin{theorem}[General paired scan loop]
\lean{Complexity.Codec.scan_loop}
\label{Complexity.Codec.scan_loop}
\leanok
\uses{Complexity.Codec.incChain, Complexity.LogProg.ARM, Complexity.LogProg.AReach, Complexity.LogProg.AReach.refl, Complexity.LogProg.AReach.single, Complexity.LogProg.AReach.step, Complexity.LogProg.AReach.trans, Complexity.LogProg.GuardLe, Complexity.LogProg.PreS, Complexity.LogProg.astep, Complexity.LogProg.register_update_le, Complexity.LogProg.vword_whole₁, Turing.pairEncode}
\difficulty{6}
Let $A$ be a register program with $m$ registers, $d$ Boolean-oracle queries, and label
type $\Lambda$. Let $y$ be its input, let $p,q,\mathrm{lim}$ be distinct registers, let
$j$ be an oracle index, and let $\delta$ be a transition function on a state and a pair
of bits. Suppose the program tests whether $p=\mathrm{lim}$, reads the first bit by
querying oracle $j$ on $y$ paired with the binary representation of $p$, reads the
second bit in the same way using $q$, applies $\delta$ to this pair, and then advances
$p$ and $q$ by increment chains of $\mathrm{stride}$ steps before returning to the test.
Assume $\mathrm{stride}>0$. Let $u,u\prime$ be equal-length bit strings, $st_0$ an
initial state, and $v$ an initial register assignment bounded by $B$. Suppose
$v(\mathrm{lim})=v(p)+\mathrm{stride}|u|$, $v(q)+\mathrm{stride}|u|\le B$, and, for each
$i<|u|$, the oracle answers at the paired inputs with binary indices
$v(p)+\mathrm{stride}i$ and $v(q)+\mathrm{stride}i$ are respectively the $i$th bits of
$u$ and $u\prime$ (taking an out-of-range bit to be false). Then the program has a
finite execution from the test label and state $st_0$ to the done label with state
obtained by folding $\delta$ over the zipped lists $u,u\prime$. The final valuation
adds $\mathrm{stride}|u|$ to registers $p$ and $q$ and leaves all other registers
unchanged; the guard holds throughout, and every register stays at most $B$.
\end{theorem}

\begin{theorem}[Stride-based decoding loop]
\lean{Complexity.Codec.decodeStride_loop}
\label{Complexity.Codec.decodeStride_loop}
\leanok
\uses{Complexity.LogProg.ARM, Complexity.LogProg.AReach, Complexity.LogProg.GuardLe, Complexity.LogProg.bitsVal, Complexity.LogProg.boundedDecodeStride_loop, Turing.pairEncode}
\difficulty{1}
Let $A$ be a register program with $m$ registers and $d$ oracle queries, with input $y$,
and let $r_{\mathrm{ptr}},r_{\mathrm{base}},r_{\mathrm{acc}},r_{\mathrm{tmp}}$ be
pairwise distinct registers. Suppose the instructions at labels
$\mathrm{test},\mathrm{downs}(i),\mathrm{read},\mathrm{copy}(b),\mathrm{bump}(b),\mathrm{clear}(b),\mathrm{scale}(b),\mathrm{take}(b),\mathrm{incs}(b,i),\mathrm{add}(b)$
have the specified loop form: test compares pointer with base; each of the first
$\mathrm{stride}$ down labels decrements the pointer and leads to read; read queries
oracle $j$ on the whole input with the pointer value; copy compares temporary with
accumulator and branches to clear or bump; bump increments temporary and returns to
copy; clear resets the accumulator; scale tests temporary for zero; take decrements it;
the first two increment labels increment the accumulator and lead back to scale; and add
either increments the accumulator and returns to test or tests temporary for zero and
returns to test. Assume these instruction equations hold with the stated continuations,
including that the decrement chain has length $\mathrm{stride}$, the increment chain has
length two, and $\mathrm{stride}>0$. Let $w$ be a bit string and $k\leq |w|$. Suppose
every suffix binary value of $w$ is at most $B$, all initial registers are at most $B$,
the pointer is initially the base value plus $\mathrm{stride}\,k$, the accumulator is
initially the binary value of the suffix of $w$ beginning at $k$, and the temporary
register is zero. Also suppose that for every $i<|w|$, the oracle answer on the pair
encoding of $y$ and the binary representation of the base value plus
$\mathrm{stride}\,i$ equals bit $i$ of $w$. Then there is a final register assignment,
still bounded by $B$, with pointer equal to its initial base value, accumulator equal to
the binary value of $w$, and temporary register zero; every other register is unchanged.
The program has a finite guarded execution from test to done, where every earlier
configuration satisfies the program's safety condition and register bound.
\end{theorem}

\begin{theorem}[Register halving loop count]
\lean{Complexity.Codec.halvCount_loop}
\label{Complexity.Codec.halvCount_loop}
\leanok
\uses{Complexity.LogProg.ARM, Complexity.LogProg.AReach, Complexity.LogProg.AReach.single, Complexity.LogProg.AReach.step, Complexity.LogProg.GuardLe, Complexity.LogProg.PreS, Complexity.LogProg.access_width_half, Complexity.LogProg.astep, Complexity.LogProg.register_update_le}
\difficulty{6}
Let $A$ be a register program with $m$ registers and an oracle, and let $y$ be an input
bit string. Choose distinct registers $t,c$ and labels $\mathrm{test},\mathrm{halve},
\mathrm{bump},\mathrm{done}$. Suppose the test label branches to $\mathrm{done}$ when
$t=0$ and otherwise to $\mathrm{halve}$; the halve label replaces $t$ by
$\lfloor t/2\rfloor$ and proceeds to $\mathrm{bump}$; and the bump label increments
$c$ and returns to $\mathrm{test}$. Let $B$ be a natural number and $v$ an initial
register assignment, with every register at most $B$. If $v(c)+|\operatorname{bits}(v(t))|
\le B$, then the program reaches $\mathrm{done}$ with $t=0$ and $c$ increased by
$|\operatorname{bits}(v(t))|$. Every configuration before the endpoint satisfies the
execution guard for bound $B$.
\end{theorem}

\begin{theorem}[Register minimum loop]
\lean{Complexity.Codec.le_loop}
\label{Complexity.Codec.le_loop}
\leanok
\uses{Complexity.LogProg.ARM, Complexity.LogProg.AReach, Complexity.LogProg.AReach.single, Complexity.LogProg.AReach.step, Complexity.LogProg.GuardLe, Complexity.LogProg.PreS, Complexity.LogProg.astep, Complexity.LogProg.register_update_le}
\difficulty{5}
Let $A$ be a register program with $m$ registers, let an oracle assign a Boolean answer
to each of $d$ indices and bit strings, and let $y$ be a bit string. Choose pairwise
distinct registers $t,a,b$ and states
$\mathrm{test}_1,\mathrm{test}_2,\mathrm{bump},\mathrm{yes},\mathrm{no}$ so that the
program branches from $\mathrm{test}_1$ to $\mathrm{yes}$ when register $t$ equals
register $a$, and otherwise to $\mathrm{test}_2$; from $\mathrm{test}_2$ it branches to
$\mathrm{no}$ when $t=b$, and otherwise to $\mathrm{bump}$; and at $\mathrm{bump}$ it
increments $t$ and returns to $\mathrm{test}_1$. If the initial register values $v$ are
all at most $B$, and $v(t)\le v(a)$ and $v(t)\le v(b)$, then the program reaches the
state selected by whether $v(a)\le v(b)$, with register $t$ replaced by
$\min(v(a),v(b))$ and every other register unchanged; throughout the run, the program
satisfies the guard requiring its register values to remain at most $B$ and its
configurations to satisfy the program's precondition.
\end{theorem}
